% Options for packages loaded elsewhere
\PassOptionsToPackage{unicode}{hyperref}
\PassOptionsToPackage{hyphens}{url}
%
\documentclass[
  ignorenonframetext,
]{beamer}
\usepackage{pgfpages}
\setbeamertemplate{caption}[numbered]
\setbeamertemplate{caption label separator}{: }
\setbeamercolor{caption name}{fg=normal text.fg}
\beamertemplatenavigationsymbolsempty
% Prevent slide breaks in the middle of a paragraph
\widowpenalties 1 10000
\raggedbottom
\setbeamertemplate{part page}{
  \centering
  \begin{beamercolorbox}[sep=16pt,center]{part title}
    \usebeamerfont{part title}\insertpart\par
  \end{beamercolorbox}
}
\setbeamertemplate{section page}{
  \centering
  \begin{beamercolorbox}[sep=12pt,center]{part title}
    \usebeamerfont{section title}\insertsection\par
  \end{beamercolorbox}
}
\setbeamertemplate{subsection page}{
  \centering
  \begin{beamercolorbox}[sep=8pt,center]{part title}
    \usebeamerfont{subsection title}\insertsubsection\par
  \end{beamercolorbox}
}
\AtBeginPart{
  \frame{\partpage}
}
\AtBeginSection{
  \ifbibliography
  \else
    \frame{\sectionpage}
  \fi
}
\AtBeginSubsection{
  \frame{\subsectionpage}
}
\usepackage{amsmath,amssymb}
\usepackage{lmodern}
\usepackage{iftex}
\ifPDFTeX
  \usepackage[T1]{fontenc}
  \usepackage[utf8]{inputenc}
  \usepackage{textcomp} % provide euro and other symbols
\else % if luatex or xetex
  \usepackage{unicode-math}
  \defaultfontfeatures{Scale=MatchLowercase}
  \defaultfontfeatures[\rmfamily]{Ligatures=TeX,Scale=1}
\fi
\usetheme[]{Montpellier}
% Use upquote if available, for straight quotes in verbatim environments
\IfFileExists{upquote.sty}{\usepackage{upquote}}{}
\IfFileExists{microtype.sty}{% use microtype if available
  \usepackage[]{microtype}
  \UseMicrotypeSet[protrusion]{basicmath} % disable protrusion for tt fonts
}{}
\makeatletter
\@ifundefined{KOMAClassName}{% if non-KOMA class
  \IfFileExists{parskip.sty}{%
    \usepackage{parskip}
  }{% else
    \setlength{\parindent}{0pt}
    \setlength{\parskip}{6pt plus 2pt minus 1pt}}
}{% if KOMA class
  \KOMAoptions{parskip=half}}
\makeatother
\usepackage{xcolor}
\IfFileExists{xurl.sty}{\usepackage{xurl}}{} % add URL line breaks if available
\IfFileExists{bookmark.sty}{\usepackage{bookmark}}{\usepackage{hyperref}}
\hypersetup{
  hidelinks,
  pdfcreator={LaTeX via pandoc}}
\urlstyle{same} % disable monospaced font for URLs
\newif\ifbibliography
\usepackage{color}
\usepackage{fancyvrb}
\newcommand{\VerbBar}{|}
\newcommand{\VERB}{\Verb[commandchars=\\\{\}]}
\DefineVerbatimEnvironment{Highlighting}{Verbatim}{commandchars=\\\{\}}
% Add ',fontsize=\small' for more characters per line
\usepackage{framed}
\definecolor{shadecolor}{RGB}{248,248,248}
\newenvironment{Shaded}{\begin{snugshade}}{\end{snugshade}}
\newcommand{\AlertTok}[1]{\textcolor[rgb]{0.94,0.16,0.16}{#1}}
\newcommand{\AnnotationTok}[1]{\textcolor[rgb]{0.56,0.35,0.01}{\textbf{\textit{#1}}}}
\newcommand{\AttributeTok}[1]{\textcolor[rgb]{0.77,0.63,0.00}{#1}}
\newcommand{\BaseNTok}[1]{\textcolor[rgb]{0.00,0.00,0.81}{#1}}
\newcommand{\BuiltInTok}[1]{#1}
\newcommand{\CharTok}[1]{\textcolor[rgb]{0.31,0.60,0.02}{#1}}
\newcommand{\CommentTok}[1]{\textcolor[rgb]{0.56,0.35,0.01}{\textit{#1}}}
\newcommand{\CommentVarTok}[1]{\textcolor[rgb]{0.56,0.35,0.01}{\textbf{\textit{#1}}}}
\newcommand{\ConstantTok}[1]{\textcolor[rgb]{0.00,0.00,0.00}{#1}}
\newcommand{\ControlFlowTok}[1]{\textcolor[rgb]{0.13,0.29,0.53}{\textbf{#1}}}
\newcommand{\DataTypeTok}[1]{\textcolor[rgb]{0.13,0.29,0.53}{#1}}
\newcommand{\DecValTok}[1]{\textcolor[rgb]{0.00,0.00,0.81}{#1}}
\newcommand{\DocumentationTok}[1]{\textcolor[rgb]{0.56,0.35,0.01}{\textbf{\textit{#1}}}}
\newcommand{\ErrorTok}[1]{\textcolor[rgb]{0.64,0.00,0.00}{\textbf{#1}}}
\newcommand{\ExtensionTok}[1]{#1}
\newcommand{\FloatTok}[1]{\textcolor[rgb]{0.00,0.00,0.81}{#1}}
\newcommand{\FunctionTok}[1]{\textcolor[rgb]{0.00,0.00,0.00}{#1}}
\newcommand{\ImportTok}[1]{#1}
\newcommand{\InformationTok}[1]{\textcolor[rgb]{0.56,0.35,0.01}{\textbf{\textit{#1}}}}
\newcommand{\KeywordTok}[1]{\textcolor[rgb]{0.13,0.29,0.53}{\textbf{#1}}}
\newcommand{\NormalTok}[1]{#1}
\newcommand{\OperatorTok}[1]{\textcolor[rgb]{0.81,0.36,0.00}{\textbf{#1}}}
\newcommand{\OtherTok}[1]{\textcolor[rgb]{0.56,0.35,0.01}{#1}}
\newcommand{\PreprocessorTok}[1]{\textcolor[rgb]{0.56,0.35,0.01}{\textit{#1}}}
\newcommand{\RegionMarkerTok}[1]{#1}
\newcommand{\SpecialCharTok}[1]{\textcolor[rgb]{0.00,0.00,0.00}{#1}}
\newcommand{\SpecialStringTok}[1]{\textcolor[rgb]{0.31,0.60,0.02}{#1}}
\newcommand{\StringTok}[1]{\textcolor[rgb]{0.31,0.60,0.02}{#1}}
\newcommand{\VariableTok}[1]{\textcolor[rgb]{0.00,0.00,0.00}{#1}}
\newcommand{\VerbatimStringTok}[1]{\textcolor[rgb]{0.31,0.60,0.02}{#1}}
\newcommand{\WarningTok}[1]{\textcolor[rgb]{0.56,0.35,0.01}{\textbf{\textit{#1}}}}
\setlength{\emergencystretch}{3em} % prevent overfull lines
\providecommand{\tightlist}{%
  \setlength{\itemsep}{0pt}\setlength{\parskip}{0pt}}
\setcounter{secnumdepth}{-\maxdimen} % remove section numbering
%\documentclass[unicode,10pt]{beamer}% 'unicode'が必要
\usepackage{luatexja}% 日本語を使う場合は必須
%\usepackage[japanese]{babel}	
%\usefonttheme[onlymath]{serif}
%\usepackage[T1]{fontenc} %これを使うと、ギリシャ文字が出ない
%\usepackage{textcomp}
%\usepackage[scale=1.0]{tgheros} % Sans serif font
%\usepackage[scaled]{beramono} % Monospace font
%\usepackage{luatexja-otf}
%\renewcommand{\kanjifamilydefault}{\gtdefault}
%\usepackage[haranoaji]{luatexja-preset}

% スライド番号の表示 %%%%%%%%%%%%%%%%%%%
\makeatletter
\setbeamertemplate{footline}{%
  \leavevmode%
  \hbox{%
  \begin{beamercolorbox}[wd=.5\paperwidth,ht=2.25ex,dp=1ex,right]{date in head/foot}%
    \insertframenumber{} \hspace*{2ex} % new version without total frames
  \end{beamercolorbox}}%
  \vskip0pt%
}
\makeatother
% スライド番号の表示 %%%%%%%%%%%%%%%%%%%
\ifLuaTeX
  \usepackage{selnolig}  % disable illegal ligatures
\fi

\title{第5章: 発見\\
(5.1. テキストデータ)}
\subtitle{今井耕介 著\\
『社会科学のためのデータ分析入門 (QSS)』}
\author{}
\date{\vspace{-2.5em}2026-03-09}

\begin{document}
\frame{\titlepage}

\hypertarget{ux30c6ux30adux30b9ux30c8ux30c7ux30fcux30bf}{%
\section{5.1
テキストデータ}\label{ux30c6ux30adux30b9ux30c8ux30c7ux30fcux30bf}}

\begin{frame}{Federalist Papers の著者推定}
\protect\hypertarget{federalist-papers-ux306eux8457ux8005ux63a8ux5b9a}{}
\begin{itemize}
\tightlist
\item
  \textbf{Federalist Papers}:
  米国憲法の批准を促すために書かれた85編のエッセイ。
\item
  ハミルトン、マディソン、ジェイの3人が執筆したが、一部の論文については著者が不明確であった。
\item
  \textbf{目標}:
  テキスト分析を用いて、著者不明の論文の執筆者を推定する。
\end{itemize}
\end{frame}

\hypertarget{ux30c6ux30adux30b9ux30c8ux30c7ux30fcux30bfux306eux51e6ux7406}{%
\section{5.1.1
テキストデータの処理}\label{ux30c6ux30adux30b9ux30c8ux30c7ux30fcux30bfux306eux51e6ux7406}}

\begin{frame}[fragile]{tm パッケージによる前処理}
\protect\hypertarget{tm-ux30d1ux30c3ux30b1ux30fcux30b8ux306bux3088ux308bux524dux51e6ux7406}{}
\begin{itemize}
\tightlist
\item
  Rの \texttt{tm}
  パッケージを使用して、文書集合（コーパス）を作成し、クリーニングを行う。
\end{itemize}

\scriptsize

\begin{Shaded}
\begin{Highlighting}[]
\FunctionTok{library}\NormalTok{(tm)}
\NormalTok{corpus.raw }\OtherTok{\textless{}{-}} \FunctionTok{VCorpus}\NormalTok{(}\FunctionTok{DirSource}\NormalTok{(}\AttributeTok{directory =} \StringTok{"federalist"}\NormalTok{, }\AttributeTok{pattern =} \StringTok{"fp"}\NormalTok{))}
\NormalTok{corpus.prep }\OtherTok{\textless{}{-}} \FunctionTok{tm\_map}\NormalTok{(corpus.raw, }\FunctionTok{content\_transformer}\NormalTok{(tolower))}
\NormalTok{corpus.prep }\OtherTok{\textless{}{-}} \FunctionTok{tm\_map}\NormalTok{(corpus.prep, stripWhitespace)}
\NormalTok{corpus.prep }\OtherTok{\textless{}{-}} \FunctionTok{tm\_map}\NormalTok{(corpus.prep, removePunctuation)}
\NormalTok{corpus.prep }\OtherTok{\textless{}{-}} \FunctionTok{tm\_map}\NormalTok{(corpus.prep, removeNumbers)}

\CommentTok{\# ステミング前の DTM (著者推定用)}
\NormalTok{dtm.pre }\OtherTok{\textless{}{-}} \FunctionTok{DocumentTermMatrix}\NormalTok{(corpus.prep)}

\CommentTok{\# ステミングありの DTM (トピック発見用)}
\NormalTok{corpus }\OtherTok{\textless{}{-}} \FunctionTok{tm\_map}\NormalTok{(corpus.prep, removeWords, }\FunctionTok{stopwords}\NormalTok{(}\StringTok{"english"}\NormalTok{))}
\NormalTok{corpus }\OtherTok{\textless{}{-}} \FunctionTok{tm\_map}\NormalTok{(corpus, stemDocument)}
\NormalTok{dtm }\OtherTok{\textless{}{-}} \FunctionTok{DocumentTermMatrix}\NormalTok{(corpus)}
\end{Highlighting}
\end{Shaded}

\normalsize
\end{frame}

\hypertarget{ux6587ux66f8ux5358ux8a9eux884cux5217-dtm}{%
\section{5.1.2 文書単語行列
(DTM)}\label{ux6587ux66f8ux5358ux8a9eux884cux5217-dtm}}

\begin{frame}[fragile]{DTM の作成}
\protect\hypertarget{dtm-ux306eux4f5cux6210}{}
\begin{itemize}
\tightlist
\item
  文書を「単語の出現頻度のベクトル」として表現する。
\end{itemize}

\scriptsize

\begin{Shaded}
\begin{Highlighting}[]
\CommentTok{\# 行列のサイズを確認}
\FunctionTok{dim}\NormalTok{(dtm)}
\end{Highlighting}
\end{Shaded}

\begin{verbatim}
## [1]   85 4849
\end{verbatim}

\begin{Shaded}
\begin{Highlighting}[]
\CommentTok{\# 一部を表示}
\FunctionTok{inspect}\NormalTok{(dtm[}\DecValTok{1}\SpecialCharTok{:}\DecValTok{5}\NormalTok{, }\DecValTok{1}\SpecialCharTok{:}\DecValTok{5}\NormalTok{])}
\end{Highlighting}
\end{Shaded}

\begin{verbatim}
## <<DocumentTermMatrix (documents: 5, terms: 5)>>
## Non-/sparse entries: 0/25
## Sparsity           : 100%
## Maximal term length: 7
## Weighting          : term frequency (tf)
## Sample             :
##           Terms
## Docs       abandon abat abb abet abhorr
##   fp01.txt       0    0   0    0      0
##   fp02.txt       0    0   0    0      0
##   fp03.txt       0    0   0    0      0
##   fp04.txt       0    0   0    0      0
##   fp05.txt       0    0   0    0      0
\end{verbatim}

\normalsize
\end{frame}

\hypertarget{ux30c8ux30d4ux30c3ux30afux306eux767aux898b}{%
\section{5.1.3
トピックの発見}\label{ux30c8ux30d4ux30c3ux30afux306eux767aux898b}}

\begin{frame}[fragile]{ワードクラウドによる可視化}
\protect\hypertarget{ux30efux30fcux30c9ux30afux30e9ux30a6ux30c9ux306bux3088ux308bux53efux8996ux5316}{}
\scriptsize

\begin{Shaded}
\begin{Highlighting}[]
\FunctionTok{library}\NormalTok{(wordcloud)}
\NormalTok{dtm.mat }\OtherTok{\textless{}{-}} \FunctionTok{as.matrix}\NormalTok{(dtm)}
\CommentTok{\# 第12論文の頻出単語}
\FunctionTok{wordcloud}\NormalTok{(}\FunctionTok{colnames}\NormalTok{(dtm.mat), dtm.mat[}\DecValTok{12}\NormalTok{, ], }\AttributeTok{max.words =} \DecValTok{20}\NormalTok{)}
\end{Highlighting}
\end{Shaded}

\begin{center}\includegraphics[height=0.6\textheight]{QSS_Ch05_1_files/figure-beamer/unnamed-chunk-3-1} \end{center}
\normalsize
\end{frame}

\begin{frame}[fragile]{tf-idf による重要語の抽出}
\protect\hypertarget{tf-idf-ux306bux3088ux308bux91cdux8981ux8a9eux306eux62bdux51fa}{}
\begin{itemize}
\tightlist
\item
  \textbf{tf-idf}:
  特定の文書で頻出するが、他の文書ではあまり使われない単語を重視する指標。
\end{itemize}

\scriptsize

\begin{Shaded}
\begin{Highlighting}[]
\NormalTok{dtm.tfidf }\OtherTok{\textless{}{-}} \FunctionTok{weightTfIdf}\NormalTok{(dtm)}
\NormalTok{dtm.tfidf.mat }\OtherTok{\textless{}{-}} \FunctionTok{as.matrix}\NormalTok{(dtm.tfidf)}
\CommentTok{\# 第12論文の重要語トップ10}
\FunctionTok{head}\NormalTok{(}\FunctionTok{sort}\NormalTok{(dtm.tfidf.mat[}\DecValTok{12}\NormalTok{, ], }\AttributeTok{decreasing =} \ConstantTok{TRUE}\NormalTok{), }\AttributeTok{n =} \DecValTok{10}\NormalTok{)}
\end{Highlighting}
\end{Shaded}

\begin{verbatim}
##     revenu contraband     patrol      excis      coast      trade        per 
## 0.01905877 0.01886965 0.01886965 0.01876560 0.01592559 0.01473504 0.01420342 
##        tax       cent     gallon 
## 0.01295466 0.01257977 0.01257977
\end{verbatim}

\normalsize
\end{frame}

\hypertarget{ux8457ux8005ux63a8ux5b9aux306eux4e88ux6e2c}{%
\section{5.1.4
著者推定の予測}\label{ux8457ux8005ux63a8ux5b9aux306eux4e88ux6e2c}}

\begin{frame}[fragile]{著者の特徴語}
\protect\hypertarget{ux8457ux8005ux306eux7279ux5fb4ux8a9e}{}
\begin{itemize}
\tightlist
\item
  ハミルトンとマディソンで使い方が異なる「機能語」（upon, there
  等）に注目する。
\end{itemize}

\scriptsize

\begin{Shaded}
\begin{Highlighting}[]
\CommentTok{\# ステミング前の行列を使用}
\NormalTok{dtm.pre.mat }\OtherTok{\textless{}{-}} \FunctionTok{as.matrix}\NormalTok{(dtm.pre)}
\CommentTok{\# 出現頻度（1000語あたり）に変換}
\NormalTok{tfm }\OtherTok{\textless{}{-}}\NormalTok{ dtm.pre.mat }\SpecialCharTok{/} \FunctionTok{rowSums}\NormalTok{(dtm.pre.mat) }\SpecialCharTok{*} \DecValTok{1000}
\NormalTok{words }\OtherTok{\textless{}{-}} \FunctionTok{c}\NormalTok{(}\StringTok{"although"}\NormalTok{, }\StringTok{"always"}\NormalTok{, }\StringTok{"commonly"}\NormalTok{, }\StringTok{"consequently"}\NormalTok{, }
           \StringTok{"considerable"}\NormalTok{, }\StringTok{"enough"}\NormalTok{, }\StringTok{"there"}\NormalTok{, }\StringTok{"upon"}\NormalTok{, }\StringTok{"while"}\NormalTok{, }\StringTok{"whilst"}\NormalTok{)}
\NormalTok{tfm }\OtherTok{\textless{}{-}}\NormalTok{ tfm[, words]}

\CommentTok{\# 既知の著者のラベル付け (Hamilton=1, Madison={-}1)}
\NormalTok{hamilton }\OtherTok{\textless{}{-}} \FunctionTok{c}\NormalTok{(}\DecValTok{1}\NormalTok{, }\DecValTok{6}\SpecialCharTok{:}\DecValTok{9}\NormalTok{, }\DecValTok{11}\SpecialCharTok{:}\DecValTok{13}\NormalTok{, }\DecValTok{15}\SpecialCharTok{:}\DecValTok{17}\NormalTok{, }\DecValTok{21}\SpecialCharTok{:}\DecValTok{36}\NormalTok{, }\DecValTok{59}\SpecialCharTok{:}\DecValTok{61}\NormalTok{, }\DecValTok{65}\SpecialCharTok{:}\DecValTok{85}\NormalTok{)}
\NormalTok{madison }\OtherTok{\textless{}{-}} \FunctionTok{c}\NormalTok{(}\DecValTok{10}\NormalTok{, }\DecValTok{14}\NormalTok{, }\DecValTok{37}\SpecialCharTok{:}\DecValTok{48}\NormalTok{, }\DecValTok{58}\NormalTok{)}
\end{Highlighting}
\end{Shaded}

\normalsize
\end{frame}

\begin{frame}[fragile]{回帰による著者予測}
\protect\hypertarget{ux56deux5e30ux306bux3088ux308bux8457ux8005ux4e88ux6e2c}{}
\scriptsize

\begin{Shaded}
\begin{Highlighting}[]
\NormalTok{author }\OtherTok{\textless{}{-}} \FunctionTok{rep}\NormalTok{(}\ConstantTok{NA}\NormalTok{, }\FunctionTok{nrow}\NormalTok{(tfm))}
\NormalTok{author[hamilton] }\OtherTok{\textless{}{-}} \DecValTok{1}\NormalTok{; author[madison] }\OtherTok{\textless{}{-}} \SpecialCharTok{{-}}\DecValTok{1}
\NormalTok{author.data }\OtherTok{\textless{}{-}} \FunctionTok{data.frame}\NormalTok{(}\AttributeTok{author =}\NormalTok{ author[}\FunctionTok{c}\NormalTok{(hamilton, madison)], }
\NormalTok{                          tfm[}\FunctionTok{c}\NormalTok{(hamilton, madison), ])}

\CommentTok{\# 4つの単語を用いた線形回帰モデル}
\NormalTok{hm.fit }\OtherTok{\textless{}{-}} \FunctionTok{lm}\NormalTok{(author }\SpecialCharTok{\textasciitilde{}}\NormalTok{ upon }\SpecialCharTok{+}\NormalTok{ there }\SpecialCharTok{+}\NormalTok{ consequently }\SpecialCharTok{+}\NormalTok{ whilst, }
             \AttributeTok{data =}\NormalTok{ author.data)}
\CommentTok{\# 著者不明の論文 (disputed) の予測}
\NormalTok{disputed }\OtherTok{\textless{}{-}} \FunctionTok{c}\NormalTok{(}\DecValTok{49}\NormalTok{, }\DecValTok{50}\SpecialCharTok{:}\DecValTok{57}\NormalTok{, }\DecValTok{62}\NormalTok{, }\DecValTok{63}\NormalTok{)}
\NormalTok{pred }\OtherTok{\textless{}{-}} \FunctionTok{predict}\NormalTok{(hm.fit, }\AttributeTok{newdata =} \FunctionTok{as.data.frame}\NormalTok{(tfm[disputed, ]))}
\NormalTok{pred }\CommentTok{\# 全て負の値であればマディソンの可能性が高い}
\end{Highlighting}
\end{Shaded}

\begin{verbatim}
##    fp49.txt    fp50.txt    fp51.txt    fp52.txt    fp53.txt    fp54.txt 
## -0.99831799 -0.06759254 -1.53243206 -0.26288400 -0.54584900 -0.56566555 
##    fp55.txt    fp56.txt    fp57.txt    fp62.txt    fp63.txt 
##  0.04376632 -0.57115610 -1.22289415 -1.00675456 -0.21939646
\end{verbatim}

\normalsize
\end{frame}

\hypertarget{ux4ea4ux5deeux59a5ux5f53ux6027-cross-validation}{%
\section{5.1.5 交差妥当性
(Cross-Validation)}\label{ux4ea4ux5deeux59a5ux5f53ux6027-cross-validation}}

\begin{frame}[fragile]{LOOCV の実行}
\protect\hypertarget{loocv-ux306eux5b9fux884c}{}
\begin{itemize}
\tightlist
\item
  1つのデータを除いてモデルを学習し、除いた1つを予測することを繰り返して、モデルの汎化性能を評価する。
\end{itemize}

\scriptsize

\begin{Shaded}
\begin{Highlighting}[]
\NormalTok{n }\OtherTok{\textless{}{-}} \FunctionTok{nrow}\NormalTok{(author.data); hm.classify }\OtherTok{\textless{}{-}} \FunctionTok{rep}\NormalTok{(}\ConstantTok{NA}\NormalTok{, n)}
\ControlFlowTok{for}\NormalTok{ (i }\ControlFlowTok{in} \DecValTok{1}\SpecialCharTok{:}\NormalTok{n) \{}
\NormalTok{    sub.fit }\OtherTok{\textless{}{-}} \FunctionTok{lm}\NormalTok{(author }\SpecialCharTok{\textasciitilde{}}\NormalTok{ upon }\SpecialCharTok{+}\NormalTok{ there }\SpecialCharTok{+}\NormalTok{ consequently }\SpecialCharTok{+}\NormalTok{ whilst, }
                  \AttributeTok{data =}\NormalTok{ author.data[}\SpecialCharTok{{-}}\NormalTok{i, ])}
\NormalTok{    hm.classify[i] }\OtherTok{\textless{}{-}} \FunctionTok{predict}\NormalTok{(sub.fit, }\AttributeTok{newdata =}\NormalTok{ author.data[i, ])}
\NormalTok{\}}
\CommentTok{\# 正答率}
\FunctionTok{mean}\NormalTok{(hm.classify[author.data}\SpecialCharTok{$}\NormalTok{author }\SpecialCharTok{==} \DecValTok{1}\NormalTok{] }\SpecialCharTok{\textgreater{}} \DecValTok{0}\NormalTok{) }\CommentTok{\# Hamilton}
\end{Highlighting}
\end{Shaded}

\begin{verbatim}
## [1] 1
\end{verbatim}

\begin{Shaded}
\begin{Highlighting}[]
\FunctionTok{mean}\NormalTok{(hm.classify[author.data}\SpecialCharTok{$}\NormalTok{author }\SpecialCharTok{==} \SpecialCharTok{{-}}\DecValTok{1}\NormalTok{] }\SpecialCharTok{\textless{}} \DecValTok{0}\NormalTok{) }\CommentTok{\# Madison}
\end{Highlighting}
\end{Shaded}

\begin{verbatim}
## [1] 1
\end{verbatim}

\normalsize
\end{frame}

\hypertarget{ux307eux3068ux3081}{%
\section{5.1.6 まとめ}\label{ux307eux3068ux3081}}

\begin{frame}{まとめ}
\protect\hypertarget{ux307eux3068ux3081-1}{}
\begin{itemize}
\tightlist
\item
  \textbf{テキストマイニング}:
  構造化されていない文章を数値化し、分析可能にする。
\item
  \textbf{前処理}: 小文字化、ストップワード削除、ステミングなどが重要。
\item
  \textbf{DTM / tf-idf}: 単語の重要度を重み付けする。
\item
  \textbf{著者推定}:
  単語の使用パターンの違いを利用して、未知の文書の執筆者を統計的に予測できる。
\item
  \textbf{交差妥当性}: モデルの信頼性を客観的に評価する。
\end{itemize}
\end{frame}

\end{document}
